% Бүлэг 2

\chapter{Судалгаа ба шийдлийн загвар} % Бүлгийн нэр
\label{Chapter2} % Энэ бүлэг рүү ишлэл хийх бол \ref{Chapter1} командыг ашигла 
%\usepackage{array}
%\usepackage{ragged2e}
%\usepackage{graphicx}

\newcolumntype{P}[1]{>{\RaggedRight\arraybackslash}p{#1}}

\section{Бүлгийн зорилго}

Энэхүү бүлэгт хүйсийн болон нийгмийн тэгш, хүртээмжтэй байдлыг хангахад чиглэсэн хиймэл оюун ухаант чатбот системийн судалгаа, шийдлийн загвар, технологийн сонголт, мэдлэгийн сан бүрдүүлэх арга зүй, хамгаалалтын зарчим болон үнэлгээний аргачлалыг тодорхойлно. Мөн ижил төстэй системүүдийг харьцуулан судалж, энэхүү дипломын ажилд RAG суурьтай архитектур сонгосон үндэслэлийг тайлбарлана.

%-------------------------------------------------------------------------------

\section{Ижил төстэй системүүдийн судалгаа ба харьцуулалт}

Ижил төстэй системүүдийг судлах нь дараах ач холбогдолтой:
\begin{itemize}
    \item бодит хэрэглээнд ашиглагдаж буй чатбот шийдлүүдийн чиг хандлагыг тодорхойлох,
    \item тэдгээрийн давуу болон сул талыг харьцуулах,
    \item өөрийн системийн шаардлага, архитектурын сонголтыг үндэслэх.
\end{itemize}

Сүүлийн жилүүдэд байгууллага, төрийн үйлчилгээний мэдээллийн системүүдэд чатбот технологи өргөн ашиглагдах болсон. Тухайлбал байгууллагын хүний нөөцийн үйлчилгээнд виртуал агент ашиглах, байгууллагын мэдлэгийн сан дээр суурилсан генератив агент хөгжүүлэх, олон нийтэд зориулсан бодлого тайлбарлах чат систем турших зэрэг жишээнүүд бий \cite{servicenow_hr_va, ms_copilot_rag, govuk_chat_guidance}.

\subsection{Харьцуулалтын шалгуур}

Ижил төстэй системүүдийг дараах шалгуураар харьцуулсан болно:
\begin{itemize}
    \item \textbf{Домэйн ба хэрэглэгч:} байгууллагын дотоод хэрэглээ, олон нийтэд зориулсан үйлчилгээ, эсвэл хосолсон хэрэглээ.
    \item \textbf{Мэдлэгийн эх сурвалж:} бодлого, журам, гарын авлага, хууль эрх зүйн эх сурвалж, FAQ зэрэгт суурилдаг эсэх.
    \item \textbf{Хариултын механизм:} дүрэмд суурилсан, retrieval-д суурилсан, эсвэл RAG/LLM ашигладаг эсэх.
    \item \textbf{Эх сурвалжийн ил тод байдал:} хариултад ашигласан эх сурвалжийг харуулах боломж.
    \item \textbf{Аюулгүй байдал:} эмзэг агуулгатай сэдэвт хязгаарлалт, чиглүүлэлт, хүний оролцоотой шийдэлтэй эсэх.
    
\end{itemize}


\begin{table}[H]
\centering
\begin{tabular}{|p{2.6cm}|p{2.0cm}|p{1.8cm}|p{1.7cm}|p{4.6cm}|}
\hline
\textbf{Систем} & 
\textbf{Домэйн} & 
\shortstack{\textbf{Архитек-}\\\textbf{тур}} & 
\shortstack{\textbf{Эх}\\\textbf{сурвалж}} & 
\textbf{Хүчтэй ба сул талууд} \\ \hline
ServiceNow HR Virtual Agent~\cite{servicenow_hr_va} &  Байгуулла- гын HR & Rule-based + retrieval & Хязгаар-лагдмал & (+) HR үйл явцад ашиглагддаг, олон хэл дэмждэг. ($-$) Лицензийн өртөг өндөр, хүйсийн тэгш байдлын тусгай домэйн мэдлэггүй \\ \hline
Microsoft Copilot Studio + RAG~\cite{ms_copilot_rag} &  Байгуулла- 
гын мэдлэгийн сан & RAG (Azure AI) & Тийм, citation-той & (+) Мэдээллийн эх сурвалжтай холбон хариулдаг, ишлэл харуулдаг. ($-$) Azure орчноос хамааралтай, Монгол хэлний дэмжлэг хязгаарлагдмал \\ \hline
UK.gov chatbot guidance~\cite{govuk_chat_guidance} & Олон нийт төрийн үйлчилгээ & Hybrid &  Бичигд- сэн & (+) Хэрэглэгч төвтэй UX-ийн зөвлөмж өгдөг. ($-$) Бодит чатбот биш, зөвхөн удирдамжийн шинжтэй \\ \hline
ChatGPT (web)~\cite{instruct_gpt,gpt3} & Универсал & Pure LLM & Үгүй (анхдагч) & (+) Уян хатан, олон хэлээр ажиллах боломжтой. ($-$) Эх сурвалжийг тогтмол харуулдаггүй, галлюцинацийн эрсдэлтэй \\ \hline
UN Women Q\&A портал~\cite{gender_equality_general} & Хүйсийн тэгш байдал & Static FAQ & Холбоос & (+) Албан ёсны эх сурвалжтай. ($-$) Чатбот биш, харилцан үйлчлэл хязгаарлагдмал \\ \hline
UNESCO bullying платформ~\cite{school_bullying_unesco} & Сургууль, насанд хүрэгчид & Static content & Тайлан, статистик & (+) Судалгаанд тулгуурласан агуулгатай. ($-$) Динамик асуулт-хариултын боломжгүй \\ \hline
{\fontsize{9}{7}\selectfont ILO HELPDESK}~\cite{ilo_c111} & Хөдөлмөр-ийн эрх & Retrieval-based & Тийм (конвенцид) & (+) Хууль, хөдөлмөрийн эрхийн эх сурвалжтай. ($-$) Интерфэйс урт, олон нийтэд ашиглахад төвөгтэй \\ \hline
\textbf{Тэгшбот} & \textbf{Хүйсийн тэгш хүртээмжтэй байдал (монгол)} & \textbf{\fontsize{9}{7}\selectfont RAG (Gemini + ChromaDB)} & \textbf{Тийм, мета өгөгдөлтэй} & \textbf{(+) Монгол хэл, Монголын хууль ба статистикт суурилсан, ишлэлтэй, crisis-routing-той} \\ \hline
\end{tabular}
\caption[Ижил төстэй системүүдийн харьцуулалт]{Ижил төстэй системүүдийн өргөтгөсөн харьцуулалт}
\label{tab:similar_systems}
\end{table}


\subsection{Харьцуулалтын хүснэгт}

Хүснэгт~\ref{tab:similar_systems}-д ижил төстэй системүүдийн харьцуулалтыг өргөтгөсөн хувилбараар үзүүлэв. Сонгосон системүүд нь олон улсын болон Монгол Улсын хэмжээнд бодит хэрэглээтэй, RAG/LLM-д суурилсан эсвэл retrieval/rule-based загвартай чатбот шийдлүүд юм.



\subsection{Судалгааны дүгнэлт}

Ижил төстэй системүүдийн харьцуулалтаас харахад орчин үеийн чатбот шийдлүүд нь хэрэглэгчийн асуултад богино хугацаанд хариулах, мэдээллийг ойлгомжтой хэлбэрээр хүргэх, байгууллагын мэдлэгийн сан болон олон нийтэд зориулсан мэдээллийн эх сурвалжтай холбогдох зэрэг давуу талтай байна. Ялангуяа RAG болон retrieval-д суурилсан системүүд нь хэрэглэгчийн асуултад хариулахдаа баримт бичиг, бодлого, журам, гарын авлага зэрэг эх сурвалжийг ашиглах боломжтойгоороо уламжлалт FAQ болон статик мэдээллийн сангаас давуу юм.

Гэсэн хэдий ч судлагдсан системүүдэд хэд хэдэн нийтлэг хязгаарлалт ажиглагдав. Үүнд Монгол хэлний дэмжлэг хязгаарлагдмал байх, хариултад ашигласан эх сурвалжийг тогтмол харуулахгүй байх, хүйсийн тэгш байдал, ялгаварлан гадуурхалт, дарамт, хүртээмжтэй байдал зэрэг эмзэг сэдэвт зориулсан тусгай мэдлэгийн сан дутмаг байх зэрэг асуудлууд багтана. Мөн зарим систем нь зөвхөн ерөнхий хэрэглээний зориулалттай тул тухайн домэйны эрх зүйн болон нийгмийн онцлогийг бүрэн тусгах боломж хязгаарлагдмал байна.

Иймээс энэхүү дипломын ажилд Монгол хэл дээр ажиллах, баримт бичигт тулгуурлан хариу үүсгэх, эх сурвалжийг ил тод харуулах боломжтой RAG суурьтай чатботын архитектурыг сонгосон. Энэ нь хэрэглэгчийн асуултад зөвхөн ерөнхий хэлний загварын мэдлэгээр бус, харин хүйсийн тэгш байдал, нийгмийн хүртээмж, ялгаварлан гадуурхалттай холбоотой найдвартай эх сурвалжид тулгуурлан хариулах нөхцөлийг бүрдүүлнэ. Мөн эмзэг агуулгатай асуултыг таних, хэрэглэгчийг тохирох мэдээлэл болон тусламжийн сувгууд руу чиглүүлэх боломжийг системийн шаардлагад тусгах үндэслэл болж байна.


%-------------------------------------------------------------------------------

\section{Системийн хамрах хүрээ}

Энэхүү чатбот системийн агуулгын хамрах хүрээнд хүйс, нас, хөгжлийн бэрхшээл, үндэс угсаа, хэл, төрсөн газар, нийгмийн гарал, эдийн засгийн байдал, бэлгийн чиг баримжаа, хүйсийн баримжаа зэрэг хүчин зүйлээс шалтгаалсан ялгаварлан гадуурхалт, гадуурхал, дээрэлхэлт, тэгш бус хандлагатай холбоотой мэдээллийг хамруулна. Мөн өсвөр үеийн орчин дахь дээрэлхэлт, боловсролын орчин дахь гадуурхал, ажлын байрны тэгш бус хандлага, байгууллагын дотоод харилцааны шударга байдал, хүртээмжтэй орчин бүрдүүлэхтэй холбоотой тайлбар, зөвлөмж, эх сурвалжийг хамруулна.

Системийн зорилтот хэрэглэгч нь өсвөр насны хэрэглэгч, оюутан залуус, ажилтан, ажил олгогч, байгууллагын хүний нөөцийн ажилтан болон эрх тэгш байдал, хүртээмжтэй орчны талаар мэдээлэл авах хүсэлтэй ерөнхий хэрэглэгчид байна. Энэ хүрээнд систем нь зөвлөгөө өгөхөөс илүүтэйгээр ойлголт тайлбарлах, эх сурвалжтай мэдээлэл хүргэх, хэрэглэгчийг зөв чиглэлд чиглүүлэх үүрэгтэй байна.

%-------------------------------------------------------------------------------



%-------------------------------------------------------------------------------


%-------------------------------------------------------------------------------

\section{Системийн шаардлагын тодорхойлолт}

Энэхүү хэсэгт системийн функциональ болон функциональ бус шаардлагыг тодорхойлно. Эдгээр шаардлага нь 3-р бүлэг дэх хэрэгжилтийн хүрээг тодорхойлж, туршилт үнэлгээнд хэмжигдэхүйц шалгуур болно.

\subsection{Функциональ шаардлага}

\begin{enumerate}
    \item Хэрэглэгчийн асуултыг хүлээн авч, текстэн хэлбэрээр хариулт өгөх чат интерфэйстэй байх.
    \item Хэрэглэгчийн асуулттай холбоотой баримтын хэсгийг мэдлэгийн сангаас хайж татах боломжтой байх.
    \item Татан авсан эх сурвалж дээр үндэслэн хариулт үүсгэх боломжтой байх.
    \item Хариултад ашигласан эх сурвалжийн хэсгийг харуулах боломжтой байх.
    \item Хэрэглэгчийн санал хүсэлт, үнэлгээг бүртгэх боломжтой байх.
    \item Системийн хэрэглээний лог, нийтлэг асуулт, сайжруулалтын мэдээллийг өгөгдлийн санд бүртгэх боломжтой байх.
    \item Эмзэг агуулгатай асуултад хязгаарлалттай, болгоомжтой, чиглүүлэлттэй хариулах боломжтой байх.
\end{enumerate}

\subsection{Функциональ бус шаардлага}

\begin{itemize}
    \item \textbf{Найдвартай байдал:} баримтад нийцсэн, ойлгомжтой хариулт өгөх.
    \item \textbf{Гүйцэтгэл:} асуултад боломжийн хугацаанд хариулах.
    \item \textbf{Нууцлал ба аюулгүй байдал:} хэрэглэгчийн мэдээллийг зохих түвшинд хамгаалах.
    \item \textbf{Хүртээмжтэй байдал:} ялгаварлан гадуурхахгүй, төвийг сахисан, хэрэглэгчид ээлтэй найруулга ашиглах.
    \item \textbf{Өргөтгөх боломж:} шинэ сэдэв, шинэ эх сурвалж, шинэ орчинд мэдлэгийн санг өргөтгөх боломжтой байх.
\end{itemize}

%-------------------------------------------------------------------------------

\section{Мэдлэгийн сан бүрдүүлэх судалгаа}

RAG системийн хамгийн чухал бүрэлдэхүүн бол мэдлэгийн сан юм. Энэхүү дипломын ажлын хүрээнд мэдлэгийн сан нь хүйсийн болон нийгмийн тэгш байдал, хүртээмжтэй байдлын талаарх төрөл бүрийн эх сурвалжаас бүрдэнэ. Үүнд:
\begin{itemize}
    \item тэгш байдал, ялгаварлан гадуурхалтаас сэргийлэх тухай байгууллагын дотоод бодлого,
    \item өсвөр үе, боловсролын орчин дахь дээрэлхэлт, гадуурхлын талаарх мэдээллийн материал,
    \item Монгол Улсын хууль эрх зүйн эх сурвалж,
    \item олон улсын байгууллагын зөвлөмж, конвенц,
    \item тайлбар, FAQ, гарын авлага зэрэг хэрэглэгч төвтэй контент
\end{itemize}
орно.

\subsection{Эх сурвалжийн төрөл}

\begin{itemize}
    \item \textbf{Албан ёсны хууль эрх зүйн эх сурвалж:} Хөдөлмөрийн тухай хууль, Жендэрийн эрх тэгш байдлыг хангах тухай хууль \cite{mn_labor_law, mn_gender_law}
    \item \textbf{Олон улсын стандарт, зөвлөмж:} ILO Convention No.111 болон тэгш боломж, гадуурхлын эсрэг зарчмууд \cite{ilo_c111}
    \item \textbf{Нийгмийн боловсролын эх сурвалж:} дээрэлхэлт, нийгмийн оролцоо зэрэг сэдвийн тайлбар материал \cite{school_bullying_unesco, gender_equality_general, social_inclusion_general}
    \item \textbf{Байгууллагын бодлого ба тайлбар контент:} ёс зүйн дүрэм, дарамтын эсрэг бодлого, FAQ, удирдамж
\end{itemize}

\subsection{Баримт боловсруулах үе шат}

Мэдлэгийн санд оруулах баримт бичгүүдийг дараах үе шатаар боловсруулна:
\begin{enumerate}
    \item \textbf{Цэвэрлэгээ:} илүүдэл тэмдэгт, header/footer, давхардсан мөрийг арилгах
    \item \textbf{Структурчлал:} гарчиг, дэд гарчиг, заалт, тайлбар хэсгийг ялгах
    \item \textbf{Хэсэглэл (chunking):} утгын хувьд холбоотой хэсгүүд болгон хуваах
    \item \textbf{Мета өгөгдөл хавсаргах:} эх сурвалжийн төрөл, огноо, сэдэв, гарчиг зэргийг тэмдэглэх
\end{enumerate}

%-------------------------------------------------------------------------------

\section{RAG архитектурын сонголт ба параметрийн судалгаа}

RAG аргачлалыг сонгох шалтгаан нь энэ системийн хариулт зөвхөн генератив биш, заавал эх сурвалжтай, тайлбарлах боломжтой байх шаардлагатайтай холбоотой. Lewis нарын судалгаанд retrieval-augmented generation нь knowledge-intensive NLP даалгаварт параметрик загвараас илүү үр дүнтэй болохыг харуулсан \cite{rag_survey}.

RAG архитектур нь хоёр үндсэн үе шаттай:
\begin{enumerate}
    \item \textbf{Retrieval:} хэрэглэгчийн асуулттай утгын хувьд төстэй баримтын хэсгүүдийг мэдлэгийн сангаас хайж олох
    \item \textbf{Generation:} татсан эх сурвалжийг контекст болгон ашиглаж хариулт үүсгэх
\end{enumerate}

\subsection{Хэсэглэлтийн стратеги}

Хэсэглэлтийн хэмжээ retrieval-ийн чанарт шууд нөлөөлдөг. Хэт жижиг chunk нь утгын холбоосыг алдагдуулж болох бол хэт том chunk нь илүүдэл контент агуулж, хамааралгүй мэдээлэл нэмэх эрсдэлтэй. Иймээс chunk size болон overlap параметрүүдийг туршилтаар тодорхойлно. Хүснэгт~\ref{tab:chunking}-д түгээмэл хэрэглэгддэг хэсэглэлтийн стратеги, тэдгээрийн давуу, сул тал, тохиромжтой хэрэглээний орчныг харьцуулсан.

\begin{table}[H]
\centering
\begin{tabular}{|p{2.8cm}|p{5.5cm}|p{5.5cm}|}
\hline
\textbf{Стратеги} & \textbf{Зарчим} & \textbf{Давуу / сул тал} \\ \hline
Тогтмол хэмжээтэй хэсэглэл
(fixed-size chunking) & Текстийг тогтмол тэмдэгт/токенаар хуваах (жишээ нь 500 тэмдэгт + 50 давхцал) & (+) Энгийн, тэнцвэртэй блок үүсгэнэ. ($-$) Өгүүлбэр болон утгын зааг дундуур тасрах эрсдэлтэй. \\ \hline
Давталттай хуваах арга
(recursive splitting) & Эхлээд догол мөр, дараа нь өгүүлбэр, эцэст нь үг гэх мэтээр түвшин түвшнээр хуваах. & (+) Утгын бүтцийг хадгална. ($-$) Параметрийн тохиргооноос ихээхэн . \\ \hline
Sentence-window & Гол өгүүлбэрийн өмнөх болон дараах өгүүлбэрүүдийг хамтад нь авч контекст үүсгэнэ. & (+) Нарийн утга, хам сэдвийг хадгална ($-$) Олон жижиг хэсэг үүсгэж, индексийн хэмжээг нэмэгдүүлнэ. \\ \hline
Утгад суурилсан хэсэглэл (semantic) & Embedding ашиглан утгын хувьд ойролцоо өгүүлбэрүүдийг нэг хэсэг болгон бүлэглэнэ. & (+) Хайлтын чанарыг сайжруулах боломжтой. ($-$) Тооцооллын зардал өндөр \\ \hline
Document-aware & Гарчиг, дэд гарчиг, хүснэгт, PDF outline зэрэг баримт бичгийн бүтцийг ашиглан хуваана. & (+) Албан ёсны баримтад тохиромжтой. ($-$) Баримт бичгийн формат болон бүтцээс хамааралтай. \\ \hline
\end{tabular}
\caption[Хэсэглэлтийн стратегийн харьцуулалт]{Хэсэглэлтийн (chunking) стратегийн харьцуулалт}
\label{tab:chunking}
\end{table}


Энэхүү дипломын ажилд эхний хувилбарт fixed-size + overlap стратегийг сонгосон бөгөөд цаашид recursive ба document-aware хэлбэрлүү шилжих боломжтойгоор pipeline-ийг зохион байгуулсан.

\subsection{Embedding ба вектор хайлт}

Embedding нь текстийн утгыг тоон вектор болгон дүрслэх арга бөгөөд retrieval шатны гол үндэс болдог \cite{sentence_bert, dpr}. Вектор хайлтын үед хэрэглэгчийн асуултын embedding болон баримтын хэсгүүдийн embedding-үүдийн ойролцоо байдлыг тооцоолж хамгийн хамааралтай хэсгийг сонгоно. ChromaDB \cite{chromadb} зэрэг тусгайлсан вектор сангууд нь metadata filter, persistent storage, олон collection-ийг шууд дэмждэг бөгөөд RAG системд түгээмэл ашиглагддаг.

Embedding загвар сонгох нь олон хэлт орчинд ялангуяа чухал. Хүснэгт~\ref{tab:embedding-models}-д өнөөгийн нийтлэг ашиглагдаж буй embedding загваруудын харьцуулалтыг үзүүлэв.

\begin{center}
\begin{table}[!htbp]
\begin{tabular}{|p{3.6cm}|p{1.5cm}|p{1.5cm}|p{1.5cm}|p{4.5cm}|}
\hline
\textbf{Загвар} & \textbf{ Хэм- жээс} & \textbf{ Контекст} & \textbf{Олон хэлт} & \textbf{Тэмдэглэл} \\ \hline
all-MiniLM-L6-v2~\cite{sentence_bert} & 384 & 512 & Хязгаар. & Хурдан, хөнгөн, English-чиглэсэн \\ \hline
multilingual-MiniLM-L12-v2~\cite{sentence_bert} & 384 & 512 & 50+ хэл & Олон хэлт энгийн baseline \\ \hline
multilingual-E5-large & 1024 & 512 & 100 хэл & Asymmetric retrieval-д тохиромжтой \\ \hline
\textbf{BAAI/bge-m3}~\cite{bge_m3} & 1024 & 8192 & 100+ хэл & Dense + sparse + multi-vector хосолсон, Монгол хэлэнд өндөр чанартай \\ \hline
OpenAI text-embedding-3-large & 3072 & 8191 & Олон хэлт & API үнэтэй, локалд ажилладаггүй \\ \hline
bge-reranker-v2-m3~\cite{bge_reranker} & скал. оноо & 8192 & 100+ хэл & Cross-encoder; re-ranking шатанд ашиглагдана \\ \hline
\end{tabular}
\caption[Embedding загваруудын харьцуулалт]{Embedding загваруудын харьцуулалт}
\label{tab:embedding-models}
\end{table}
\end{center}

Энэхүү судалгаанд BAAI/bge-m3 загварыг сонгосон гол шалтгаан нь Монгол Cyrillic-ийг агуулсан 100 гаруй хэлийг дэмждэг, dense ба sparse retrieval-ийг хослуулсан, 8192 token-ийн урт контекст дэмждэг, нээлттэй эх зөвшөөрөлтэйг харгалзан үзсэн юм.

\subsection{Top-$k$ ба дахин эрэмбэлэлт}

Retrieval шатанд top-$k$ гэдэг нь хамгийн өндөр хамааралтай хэдэн баримтын хэсгийг сонгохыг илэрхийлнэ. Хэрэв $k$ хэт бага байвал чухал эх сурвалж алдагдаж болзошгүй, харин хэт их байвал prompt-д илүүдэл мэдээлэл орж хариултын чанарт сөргөөр нөлөөлж болзошгүй. Иймээс top-$k$-ийн оновчтой утгыг туршилтаар тогтооно. Шаардлагатай тохиолдолд reranking ашиглан илүү хамааралтай эх сурвалжийг урьтал болгоно.

\subsection{Prompt загвар ба эх сурвалжийн ишлэл}

Хариулт үүсгэх prompt нь дараах зарчмыг баримтална:
\begin{itemize}
    \item асуултын зорилгыг тодорхойлох,
    \item retrieval-ээр татсан эх сурвалжийг контекст болгон оруулах,
    \item баримтаас хазайхгүй байх дүрэм тогтоох,
    \item боломжтой тохиолдолд эх сурвалжийн үндэслэл харуулах.
\end{itemize}

Ингэснээр grounded generation буюу эх сурвалжид суурилсан хариулт үүсгэх боломж нэмэгдэнэ \cite{grounded_generation}.

%-------------------------------------------------------------------------------

\section{Хариуцлагатай AI ба хамгаалалтын зохицуулалт}

Тэгш байдал, гадуурхал, дээрэлхэл, нийгмийн оролцоо зэрэг нь эмзэг агуулга бүхий сэдвүүд тул чатбот систем нь болгоомжтой, хариуцлагатай хариулах шаардлагатай. Иймээс дараах хамгаалалтын зарчмуудыг тусгана:
\begin{enumerate}
    \item \textbf{Ялгаварлан гадуурхахгүй найруулга:} хэвшмэл ойлголт давтахгүй, төвийг сахисан, хүндэтгэлтэй хэл ашиглах
    \item \textbf{Хэт эрсдэлтэй зөвлөгөөнөөс татгалзах:} хууль зүйн эцсийн дүгнэлт, эмчилгээний заавар зэрэг өндөр эрсдэлтэй зөвлөгөөг шууд өгөхгүй байх
    \item \textbf{Чиглүүлэлт өгөх:} хэрэглэгчийг холбогдох эх сурвалж, албан ёсны байгууллага, үйлчилгээ, бодлогын баримт бичиг рүү чиглүүлэх
    \item \textbf{Лог ба аудит:} системийн хариулт, retrieval, хэрэглэгчийн үнэлгээ зэргийг бүртгэн сайжруулалтын үндэс болгох
\end{enumerate}

Responsible AI-ийн fairness, transparency, accountability, inclusiveness зэрэг зарчим нь ийм төрлийн системийн үндсэн шаардлага гэж үздэг \cite{responsible_ai}.

%-------------------------------------------------------------------------------

\section{Туршилт ба үнэлгээний аргачлал}

Энэхүү системийн үнэлгээг гурван үндсэн чиглэлээр хийнэ: retrieval чанар, хариултын чанар, хэрэглэгчийн үнэлгээ. Хүснэгт~\ref{tab:eval-metrics}-д ашиглах гол үнэлгээний хэмжүүрүүдийг танилцуулав.


\begin{table}[H]
\centering
\begin{tabular}{|p{2.4cm}|p{4.0cm}|p{7.5cm}|}
\hline
\textbf{Шат} & \textbf{Хэмжүүр} & \textbf{Тайлбар} \\ \hline
Retrieval & Precision@$k$, Recall@$k$~\cite{dpr} & Top-$k$ дотор хэдэн зөв баримт орсон, хэдэн зөв баримт алдагдсаныг тооцно \\ \hline
Retrieval & Mean Reciprocal Rank (MRR) & Зөв баримтын дундаж урвуу эрэмбийг тооцох; эрэмбэлэлтийн чанарыг хэмжинэ \\ \hline
Retrieval & nDCG@$k$ & Эрэмбэлэгдсэн үр дүнг ranking-аар нь жигнэсэн чанарын үзүүлэлт \\ \hline
Generation & Faithfulness~\cite{ragas} & Хариулт татсан контекстээс хазайсан эсэх (RAGAS) \\ \hline
Generation & Answer Relevancy~\cite{ragas} & Хариулт асуулттай шууд холбогдсон эсэх \\ \hline
Generation & Context Precision / Recall~\cite{ragas} & Татсан контекст хариултанд хэр хэрэгцээтэй байсныг хэмжих \\ \hline
Generation & Hallucination rate~\cite{hallucination_survey,ji_hallucination_survey} & Эх сурвалжид байхгүй зохиомол мэдээллийн эзлэх хувь \\ \hline
User & Likert (1--5) үнэлгээ & Хэрэглэгчээс ойлгомжтой, итгэл төрсөн, хэрэгтэй эсэх асуултаар цуглуулна \\ \hline
User & Task Success Rate & Хэрэглэгч асуултынхаа хариулт авч чадсан эсэхийг тооцох \\ \hline
System & Latency (ms) & Асуултын хариулт буцах хүртэлх дундаж хугацаа \\ \hline
System & Token usage & LLM-ийн нэг хариулт үүсгэхэд зарцуулсан токен (зардлын үзүүлэлт) \\ \hline
\end{tabular}
\caption[RAG системийн үнэлгээний хэмжүүрүүд]{RAG системийн үнэлгээний хэмжүүрүүд}
\label{tab:eval-metrics}
\end{table}


\subsection{Retrieval чанарын үнэлгээ}

Retrieval шатанд зөв эх сурвалж top-$k$ дотор орж ирж байгаа эсэхийг Precision@$k$, Recall@$k$, MRR ба nDCG зэрэг сонгодог IR хэмжүүрээр хэмжинэ \cite{dpr}. Туршилтын асуултын багц нь хүйсийн тэгш байдал, дарамт, дээрэлхэлт, гадуурхалт, тусламжийн утас зэрэг сэдэв тус бүрд 10--15 асуулт агуулсан байх шаардлагатай. Хэсэглэлтийн хэмжээ, embedding загвар, top-$k$ зэрэг параметрийн нөлөөг харьцуулж үнэлснээр оновчтой тохиргоог тогтооно.

\subsection{Хариултын чанарын үнэлгээ}

Хариултын чанарыг дараах шалгуураар үнэлнэ:
\begin{itemize}
    \item \textbf{Нийцтэй байдал (Relevance):} хариулт асуулттай шууд холбоотой эсэх \cite{ragas}
    \item \textbf{Ойлгомжтой байдал (Clarity):} хэрэглэгч ойлгоход хялбар эсэх
    \item \textbf{Эх сурвалжтай нийцэл (Faithfulness):} хариулт татсан баримтаас хазайсан эсэх \cite{ragas, hallucination_survey}
    \item \textbf{Хэл найруулгын зохистой байдал:} ялгаварлан гадуурхахгүй, төвийг сахисан байдал
\end{itemize}

Эдгээр шалгуурыг хэмжихэд RAGAS framework \cite{ragas}-ийг ашиглах боломжтой бөгөөд тус нь LLM-ийг үнэлгээний хэрэгсэл болгон ашиглан faithfulness, answer relevancy, context precision/recall зэрэг автомат хэмжүүрүүдийг тооцдог. Хүний эксперт үнэлгээтэй хослуулснаар үнэлгээний бодит чанарыг хангах боломжтой.

\subsection{Хэрэглэгчийн үнэлгээ}

Хэрэглэгчээс 1--5 онооны үнэлгээ, мөн чанарын санал хүсэлт цуглуулна. Үүнд:
\begin{itemize}
    \item мэдээлэл ойлгомжтой байсан эсэх,
    \item хариулт хэрэгтэй байсан эсэх,
    \item эх сурвалжтай байдал итгэл төрүүлсэн эсэх
\end{itemize}
зэрэг шалгуур орно.

%-------------------------------------------------------------------------------

\section{Бүлгийн дүгнэлт}

Энэхүү бүлэгт ижил төстэй системүүдийн судалгаа, системийн хамрах хүрээ, шаардлагын тодорхойлолт, мэдлэгийн сан бүрдүүлэх арга зүй, RAG архитектурын сонголт, хэсэглэлтийн стратеги, embedding загваруудын харьцуулалт, хариуцлагатай AI-ийн хамгаалалтын зарчим болон үнэлгээний аргачлалыг тодорхойллоо. Судалгаанаас үзэхэд хүйсийн болон нийгмийн тэгш, хүртээмжтэй байдлын талаарх чатбот систем нь заавал баримтад тулгуурласан, тайлбарлах боломжтой, ялгаварлан гадуурхахгүй байх шаардлагатай байна. Embedding загвар, хэсэглэлтийн стратеги, retrieval ба generation хоёр шатны үнэлгээний хэмжүүрүүдийг тодорхой болгосон нь дараагийн бүлэгт системийг хэрэгжүүлэх, туршихад баримтлах үндэс суурь болж байна. Иймээс энэхүү дипломын ажилд RAG суурьтай (BAAI/bge-m3 embedding + ChromaDB + Gemini), эх сурвалжийн ил тод байдал бүхий чатбот архитектурыг сонгож, дараагийн бүлэгт түүний хэрэгжилтийг авч үзнэ.

